\documentclass[11pt]{amsart}
\usepackage{amssymb, amsthm}
\usepackage[shortlabels]{enumitem}
\usepackage[margin=1in]{geometry}

\theoremstyle{definition}
\newtheorem{exercise}{Exercise}
\newtheorem*{remark}{Remark}

\title{Homework 3: Sets, Logic, and Induction}
\author{Math 101 --- Introduction to Proofs}
\date{Spring 2026}

\begin{document}
\maketitle

\textbf{Instructions.} Write complete, clearly organized proofs. State which proof technique you are using for each problem (direct proof, contrapositive, contradiction, or induction).

\begin{exercise}
Write each of the following sets by listing their elements.
\begin{enumerate}[(a)]
  \item $\{ x \in \mathbb{Z} : -3 \leq x < 5 \}$
  \item $\{ x \in \mathbb{N} : x^2 < 20 \}$
  \item $\{ 2k + 1 : k \in \mathbb{Z}, \; 0 \leq k \leq 4 \}$
\end{enumerate}
\end{exercise}

\begin{exercise}
Write each of the following sets using set-builder notation.
\begin{enumerate}[(a)]
  \item $\{1, 4, 9, 16, 25, \ldots\}$
  \item $\{\ldots, -8, -4, 0, 4, 8, 12, \ldots\}$
\end{enumerate}
\end{exercise}

\begin{exercise}
Negate each of the following statements. Your negation should not begin with ``It is not the case that.''
\begin{enumerate}[(a)]
  \item For every real number $x$, if $x > 0$ then $x^2 > 0$.
  \item There exists an integer $n$ such that $n^2 + n$ is odd.
  \item For all $\varepsilon > 0$, there exists $\delta > 0$ such that $|x - a| < \delta$ implies $|f(x) - f(a)| < \varepsilon$.
\end{enumerate}
\end{exercise}

\begin{exercise}
Prove each of the following statements.
\begin{enumerate}[(a)]
  \item If $n$ is an odd integer, then $n^2$ is odd.
  \item If $a \mid b$ and $a \mid c$, then $a \mid (b + c)$.
  \item If $n^2$ is even, then $n$ is even. \emph{(Hint: Try the contrapositive.)}
  \item $\sqrt{3}$ is irrational. \emph{(Hint: Use proof by contradiction.)}
\end{enumerate}
\end{exercise}

\begin{exercise}
Prove by mathematical induction that for all positive integers $n$:
\[
  \sum_{k=1}^{n} k^2 = \frac{n(n+1)(2n+1)}{6}.
\]
\end{exercise}

\begin{exercise}
Let $A$, $B$, and $C$ be sets. Prove that $A \cap (B \cup C) = (A \cap B) \cup (A \cap C)$.
\end{exercise}

\begin{exercise}
Suppose $f \colon A \to B$ and $g \colon B \to C$ are functions.
\begin{enumerate}[(a)]
  \item Prove that if $g \circ f$ is injective, then $f$ is injective.
  \item Prove that if $g \circ f$ is surjective, then $g$ is surjective.
  \item Give an example where $g \circ f$ is bijective but neither $f$ nor $g$ is bijective. Or prove that no such example exists.
\end{enumerate}
\end{exercise}

\begin{remark}
Parts (a) and (b) of the previous exercise illustrate a general principle: composition ``inherits'' injectivity from the first function applied and surjectivity from the second.
\end{remark}

\end{document}
